\documentclass[11pt,a4paper]{article}
\usepackage[top=1.6cm,bottom=1.4cm,left=1.3cm,right=1.3cm,headsep=8pt,footskip=16pt]{geometry}
\usepackage{enumitem}
\usepackage{amssymb}
\usepackage{array}
\usepackage{xcolor}
\usepackage{fancyhdr}
\usepackage{multicol}
\usepackage{polyglossia}
\setdefaultlanguage[calendar=gregorian, numerals=maghrib]{arabic}
\setotherlanguage{english}
\IfFontExistsTF{Amiri}{%
  \newfontfamily\arabicfont[Script=Arabic, Scale=1.08]{Amiri}
}{%
  \newfontfamily\arabicfont[Script=Arabic, Scale=1.08, Path=../fonts/, Extension=.ttf,
    UprightFont=*-Regular, BoldFont=*-Bold, ItalicFont=*-Italic,
    BoldItalicFont=*-BoldItalic]{Amiri}
}
\IfFontExistsTF{AlBattar}{%
  \newfontfamily\khatt[Script=Arabic, Scale=1.15]{AlBattar}
}{%
  \newfontfamily\khatt[Script=Arabic, Scale=1.15, Path=../fonts/,
    BoldFont=ae_AlBattar.ttf]{ae_AlBattar.ttf}
}
\renewcommand{\headrulewidth}{0.8pt}
\renewcommand{\footrulewidth}{0.8pt}
\fancyhf{}
\rhead{\khatt المدرسة العليا للأساتذة ببوزريعة}
\lhead{\khatt ملحقة بومرداس}
\lfoot{\khatt الأستاذ: د. بن عبد الرزاق يوسف}
\cfoot{\thepage/2}
\pagestyle{fancy}
\parindent 0pt
\setlength{\columnsep}{0.6cm}
\setlength{\columnseprule}{0.4pt}
\setlist[enumerate,1]{label=\textbf{\arabic*.}, leftmargin=*, itemsep=5pt, topsep=2pt}
\begin{document}

\begin{center}
{\Large\khatt\textbf{الاختبار الثاني في الذكاء الاصطناعي}}\\[2pt]
{\normalsize أسئلة متعددة الاختيارات ــ الفصول من 1 إلى 6 ــ 15 سؤالًا}
\end{center}
\vspace{-4pt}\rule{\textwidth}{0.6pt}\\
{\small\textbf{التعليمة:} لكل سؤال إجابة صحيحة واحدة فقط من بين خمسة اختيارات.
اكتب \textbf{حرف} الإجابة الصحيحة في جدول الإجابات أسفل الصفحة الثانية.
لكل سؤال نقطة واحدة (المجموع 15 نقطة).}
\vspace{-2pt}
\begin{multicols}{2}

\noindent\begin{minipage}{\linewidth}
\hangindent=0.7cm\hangafter=1\makebox[0.7cm][r]{\textbf{1.}}\textbf{ماذا تكتشف الآلة في البيانات أثناء التعلّم؟}
\begin{itemize}[leftmargin=0.95cm, labelsep=4pt, itemsep=0pt, topsep=2pt, parsep=0pt]
  \item[\textbf{أ)}] الأنماط التي تتكرر
  \item[\textbf{ب)}] كلمات المرور المخزنة
  \item[\textbf{ج)}] الأخطاء الإملائية فقط
  \item[\textbf{د)}] الإعلانات التجارية
  \item[\textbf{هـ)}] أسماء كتّاب النصوص
\end{itemize}
\end{minipage}\par\vspace{7pt plus 4pt}
\noindent\begin{minipage}{\linewidth}
\hangindent=0.7cm\hangafter=1\makebox[0.7cm][r]{\textbf{2.}}\textbf{أين نضع جملة الإفصاح في البحث؟}
\begin{itemize}[leftmargin=0.95cm, labelsep=4pt, itemsep=0pt, topsep=2pt, parsep=0pt]
  \item[\textbf{أ)}] لا نضعها في أي مكان أبدًا
  \item[\textbf{ب)}] بخط مخفي داخل الفقرات
  \item[\textbf{ج)}] في اسم الملف المرسَل فقط
  \item[\textbf{د)}] في رسالة خاصة إلى الأداة
  \item[\textbf{هـ)}] في آخر البحث أو في مقدمته
\end{itemize}
\end{minipage}\par\vspace{7pt plus 4pt}
\noindent\begin{minipage}{\linewidth}
\hangindent=0.7cm\hangafter=1\makebox[0.7cm][r]{\textbf{3.}}\textbf{ما «النمط» في هذا الدرس؟}
\begin{itemize}[leftmargin=0.95cm, labelsep=4pt, itemsep=0pt, topsep=2pt, parsep=0pt]
  \item[\textbf{أ)}] سرعة الأداة في إنتاج الإجابة
  \item[\textbf{ب)}] شكل المعلومة: نص أو صورة أو صوت
  \item[\textbf{ج)}] الشركة التي طوّرت الأداة
  \item[\textbf{د)}] اللغة التي تكتب بها الأداة
  \item[\textbf{هـ)}] عدد الكلمات في الأمر المكتوب
\end{itemize}
\end{minipage}\par\vspace{7pt plus 4pt}
\noindent\begin{minipage}{\linewidth}
\hangindent=0.7cm\hangafter=1\makebox[0.7cm][r]{\textbf{4.}}\textbf{في خطوات التحقق من المعلومة، ما الخطوة التي تأتي مباشرة بعد «حدّد المعلومة»؟}
\begin{itemize}[leftmargin=0.95cm, labelsep=4pt, itemsep=0pt, topsep=2pt, parsep=0pt]
  \item[\textbf{أ)}] ابحث في مصدر موثوق
  \item[\textbf{ب)}] قرّر هل تعتمدها
  \item[\textbf{ج)}] قارن بين الإجابتين
  \item[\textbf{د)}] اكتب طلبًا جديدًا
  \item[\textbf{هـ)}] اقرأ الإجابة كاملة
\end{itemize}
\end{minipage}\par\vspace{7pt plus 4pt}
\noindent\begin{minipage}{\linewidth}
\hangindent=0.7cm\hangafter=1\makebox[0.7cm][r]{\textbf{5.}}\textbf{ماذا يقدّم النموذج في مرحلة «النتيجة»؟}
\begin{itemize}[leftmargin=0.95cm, labelsep=4pt, itemsep=0pt, topsep=2pt, parsep=0pt]
  \item[\textbf{أ)}] قواعد مكتوبة يدويًا
  \item[\textbf{ب)}] تقسيمًا جديدًا للبيانات
  \item[\textbf{ج)}] صورًا جديدة لتدريب نفسه
  \item[\textbf{د)}] توقعًا أو تصنيفًا
  \item[\textbf{هـ)}] بيانات اختبار إضافية
\end{itemize}
\end{minipage}\par\vspace{7pt plus 4pt}
\noindent\begin{minipage}{\linewidth}
\hangindent=0.7cm\hangafter=1\makebox[0.7cm][r]{\textbf{6.}}\textbf{ما الفرق بين بيانات التدريب وبيانات الاختبار؟}
\begin{itemize}[leftmargin=0.95cm, labelsep=4pt, itemsep=0pt, topsep=2pt, parsep=0pt]
  \item[\textbf{أ)}] لا فرق بينهما، فهي البيانات نفسها تمامًا
  \item[\textbf{ب)}] الأولى لتقييم النموذج، والثانية ليتعلم منها
  \item[\textbf{ج)}] الأولى يتعلم منها النموذج، والثانية لتقييمه
  \item[\textbf{د)}] الأولى صغيرة الحجم، والثانية كبيرة دائمًا
  \item[\textbf{هـ)}] الأولى صور دائمًا، والثانية نصوص دائمًا
\end{itemize}
\end{minipage}\par\vspace{7pt plus 4pt}
\noindent\begin{minipage}{\linewidth}
\hangindent=0.7cm\hangafter=1\makebox[0.7cm][r]{\textbf{7.}}\textbf{حصلت على إجابة غير مرضية من الأداة. ماذا تفعل؟}
\begin{itemize}[leftmargin=0.95cm, labelsep=4pt, itemsep=0pt, topsep=2pt, parsep=0pt]
  \item[\textbf{أ)}] تعيد الأمر نفسه دون تغيير
  \item[\textbf{ب)}] تترك الأداة نهائيًا ولا تعود
  \item[\textbf{ج)}] تحذف كل عناصر الأمر وتبدأ
  \item[\textbf{د)}] تعدّل الأمر وتعيد المحاولة
  \item[\textbf{هـ)}] تقبل أول إجابة في كل مرة
\end{itemize}
\end{minipage}\par\vspace{7pt plus 4pt}
\noindent\begin{minipage}{\linewidth}
\hangindent=0.7cm\hangafter=1\makebox[0.7cm][r]{\textbf{8.}}\textbf{طلب أفكار من الأداة، ثم كتابة الإجابة بنفسي، استعمال:}
\begin{itemize}[leftmargin=0.95cm, labelsep=4pt, itemsep=0pt, topsep=2pt, parsep=0pt]
  \item[\textbf{أ)}] انتحال
  \item[\textbf{ب)}] مرفوض
  \item[\textbf{ج)}] غير مسؤول
  \item[\textbf{د)}] ممنوع دائمًا
  \item[\textbf{هـ)}] مقبول
\end{itemize}
\end{minipage}\par\vspace{7pt plus 4pt}
\noindent\begin{minipage}{\linewidth}
\hangindent=0.7cm\hangafter=1\makebox[0.7cm][r]{\textbf{9.}}\textbf{عند تضييق الموضوع «الذكاء الاصطناعي»، ما الصياغة الأدق؟}
\begin{itemize}[leftmargin=0.95cm, labelsep=4pt, itemsep=0pt, topsep=2pt, parsep=0pt]
  \item[\textbf{أ)}] «التكنولوجيا الحديثة وتأثيرها على الحياة»
  \item[\textbf{ب)}] «فوائد الذكاء الاصطناعي في التعليم الجامعي ومخاطره»
  \item[\textbf{ج)}] «الحاسوب والإنترنت والهاتف والذكاء الاصطناعي»
  \item[\textbf{د)}] «التعليم والتكنولوجيا والمجتمع والاقتصاد»
  \item[\textbf{هـ)}] «الذكاء الاصطناعي في العالم وفي كل المجالات»
\end{itemize}
\end{minipage}\par\vspace{7pt plus 4pt}
\noindent\begin{minipage}{\linewidth}
\hangindent=0.7cm\hangafter=1\makebox[0.7cm][r]{\textbf{10.}}\textbf{رفع أستاذ مخططًا بيانيًا، فطلب شرحًا مبسطًا، ثم طلب قراءته صوتيًا. ما تسلسل الأنماط؟}
\begin{itemize}[leftmargin=0.95cm, labelsep=4pt, itemsep=0pt, topsep=2pt, parsep=0pt]
  \item[\textbf{أ)}] صوت $\leftarrow$ نص $\leftarrow$ صورة
  \item[\textbf{ب)}] صورة $\leftarrow$ نص $\leftarrow$ صوت
  \item[\textbf{ج)}] نص $\leftarrow$ صورة $\leftarrow$ صوت
  \item[\textbf{د)}] نص $\leftarrow$ صوت $\leftarrow$ صورة
  \item[\textbf{هـ)}] صورة $\leftarrow$ صوت $\leftarrow$ نص
\end{itemize}
\end{minipage}\par\vspace{7pt plus 4pt}
\noindent\begin{minipage}{\linewidth}
\hangindent=0.7cm\hangafter=1\makebox[0.7cm][r]{\textbf{11.}}\textbf{أي عبارة تلخّص بدقة مكانة الذكاء الاصطناعي في عمل الأستاذ؟}
\begin{itemize}[leftmargin=0.95cm, labelsep=4pt, itemsep=0pt, topsep=2pt, parsep=0pt]
  \item[\textbf{أ)}] مساعد يقترح، والأستاذ يختار ويبقى مسؤولًا
  \item[\textbf{ب)}] مساعد يتحمل مسؤولية الأخطاء بدل الأستاذ
  \item[\textbf{ج)}] مرجع علمي نهائي يغني عن المعاجم والدواوين
  \item[\textbf{د)}] أداة لا فائدة منها إلا في ترجمة النصوص
  \item[\textbf{هـ)}] بديل كامل عن الأستاذ في الإعداد والتقويم
\end{itemize}
\end{minipage}\par\vspace{7pt plus 4pt}
\noindent\begin{minipage}{\linewidth}
\hangindent=0.7cm\hangafter=1\makebox[0.7cm][r]{\textbf{12.}}\textbf{الأمر: «اقترح عشر جمل فيها مفعول لأجله، لتلاميذ السنة الثالثة متوسط». ما العناصر الموجودة فيه؟}
\begin{itemize}[leftmargin=0.95cm, labelsep=4pt, itemsep=0pt, topsep=2pt, parsep=0pt]
  \item[\textbf{أ)}] الدور والتنسيق والسياق
  \item[\textbf{ب)}] القيود والدور والتنسيق
  \item[\textbf{ج)}] السياق والتنسيق والدور
  \item[\textbf{د)}] الدور والمهمة والتنسيق
  \item[\textbf{هـ)}] المهمة والقيود والسياق
\end{itemize}
\end{minipage}\par\vspace{7pt plus 4pt}
\noindent\begin{minipage}{\linewidth}
\hangindent=0.7cm\hangafter=1\makebox[0.7cm][r]{\textbf{13.}}\textbf{أي عبارة تعبّر بدقة عن فكرة الدرس حول الأخلاقيات؟}
\begin{itemize}[leftmargin=0.95cm, labelsep=4pt, itemsep=0pt, topsep=2pt, parsep=0pt]
  \item[\textbf{أ)}] الأداة غير أخلاقية دائمًا، فيجب منعها في التعليم كليًا
  \item[\textbf{ب)}] الأداة أخلاقية بطبيعتها، فلا حاجة إلى أي قواعد للاستعمال
  \item[\textbf{ج)}] الأداة لا تقرر كيف تُستعمل؛ الإنسان يختار الأمانة أو الغش
  \item[\textbf{د)}] الأخلاقيات تخص الأساتذة وحدهم، لا الطلبة ولا التلاميذ
  \item[\textbf{هـ)}] الأخلاقيات تخص المطوّرين وحدهم، لا المستخدمين العاديين
\end{itemize}
\end{minipage}\par\vspace{7pt plus 4pt}
\noindent\begin{minipage}{\linewidth}
\hangindent=0.7cm\hangafter=1\makebox[0.7cm][r]{\textbf{14.}}\textbf{رتّب المصادر من الأكثر موثوقية إلى الأقل:}
\begin{itemize}[leftmargin=0.95cm, labelsep=4pt, itemsep=0pt, topsep=2pt, parsep=0pt]
  \item[\textbf{أ)}] مقال صحفي $\leftarrow$ مقال محكّم $\leftarrow$ موقع جامعة $\leftarrow$ محتوى ذكاء اصطناعي $\leftarrow$ مدونة
  \item[\textbf{ب)}] مقال محكّم $\leftarrow$ مدونة $\leftarrow$ موقع جامعة $\leftarrow$ مقال صحفي $\leftarrow$ محتوى ذكاء اصطناعي
  \item[\textbf{ج)}] محتوى ذكاء اصطناعي $\leftarrow$ مقال محكّم $\leftarrow$ موقع جامعة $\leftarrow$ مقال صحفي $\leftarrow$ مدونة
  \item[\textbf{د)}] مقال محكّم $\leftarrow$ موقع جامعة $\leftarrow$ مقال صحفي $\leftarrow$ مدونة $\leftarrow$ محتوى ذكاء اصطناعي
  \item[\textbf{هـ)}] موقع جامعة $\leftarrow$ مقال محكّم $\leftarrow$ مدونة $\leftarrow$ مقال صحفي $\leftarrow$ محتوى ذكاء اصطناعي
\end{itemize}
\end{minipage}\par\vspace{7pt plus 4pt}
\noindent\begin{minipage}{\linewidth}
\hangindent=0.7cm\hangafter=1\makebox[0.7cm][r]{\textbf{15.}}\textbf{رفع أستاذ صورة غير واضحة لمخطوطة قديمة بخط مغربي. ما التصرف الأنسب؟}
\begin{itemize}[leftmargin=0.95cm, labelsep=4pt, itemsep=0pt, topsep=2pt, parsep=0pt]
  \item[\textbf{أ)}] يحوّل الصورة إلى صوت لتصبح الحروف أوضح
  \item[\textbf{ب)}] يكتفي بقراءة شرح الأداة دون مقارنة بالأصل
  \item[\textbf{ج)}] يختار صورة أوضح، ثم يقارن النص الناتج بالمخطوطة
  \item[\textbf{د)}] يعتمد النص الناتج مباشرة دون أي مراجعة له
  \item[\textbf{هـ)}] يترك الأداة تخمّن الكلمات الغامضة ويعتمدها
\end{itemize}
\end{minipage}\par\vspace{7pt plus 4pt}
\end{multicols}
\vfill
\begin{center}
{\khatt\textbf{جدول الإجابات}}\\[3pt]
\renewcommand{\arraystretch}{1.6}
\begin{tabular}{|>{\centering\arraybackslash}m{1.9cm}|*{15}{>{\centering\arraybackslash}m{0.62cm}|}}
\hline
\textbf{رقم السؤال} & \textbf{1} & \textbf{2} & \textbf{3} & \textbf{4} & \textbf{5} & \textbf{6} & \textbf{7} & \textbf{8} & \textbf{9} & \textbf{10} & \textbf{11} & \textbf{12} & \textbf{13} & \textbf{14} & \textbf{15} \\ \hline
\textbf{الإجابة} &  &  &  &  &  &  &  &  &  &  &  &  &  &  &  \\[8pt] \hline
\end{tabular}

\vspace{8pt}
\renewcommand{\arraystretch}{2.2}
\begin{tabular}{|>{\centering\arraybackslash}m{7.5cm}|>{\centering\arraybackslash}m{3.2cm}|>{\centering\arraybackslash}m{3.2cm}|}
\hline
\textbf{الاسم واللقب} & \textbf{الفوج} & \textbf{العلامة} \\ \hline
 & & \hfill /15 \\ \hline
\end{tabular}
\end{center}
\end{document}
